\documentclass[11pt,a4paper]{article}
\usepackage[utf8]{inputenc}
\usepackage[T1]{fontenc}
\usepackage{amsmath,amsfonts,amssymb}
\usepackage{graphicx}
\usepackage{hyperref}
\usepackage{listings}
\usepackage{color}
\usepackage{booktabs}
\usepackage{float}
\usepackage{algorithm}
\usepackage{algorithmic}
\usepackage{geometry}
\geometry{margin=1in}
\definecolor{zenblue}{RGB}{41,121,255}
\hypersetup{colorlinks=true,linkcolor=zenblue,urlcolor=zenblue,citecolor=zenblue}

\title{\textbf{Zen Training Methodology: From Pretraining to Deployment}\\
\large Technical Report v2025.03}
\author{Antje Worring, Zach Kelling \\ Zen LM Research Team\\
\texttt{research@zenlm.org}}
\date{March 2025}

\begin{document}
\maketitle

\begin{abstract}
We present the complete post-training and continued-training methodology for the Zen
family of large language models, which are built on the open-weight Qwen3 models
\cite{qwen3} (Apache-2.0): dense 0.6B/4B/8B/32B variants and the Qwen3-30B-A3B
mixture-of-experts variant. Starting from the Qwen3 base checkpoints, our pipeline spans
continued pretraining on curated data through supervised fine-tuning and reinforcement
learning from human feedback. Our approach introduces several key components: a hybrid
data mixing strategy that balances quality and diversity, perplexity-based quality filtering
that removes low-signal content, and a constitutional training regime that instills
safety properties without sacrificing capability. We report training curves, compute
budget analysis, and ablation results demonstrating the contribution of each
methodological component. The same pipeline applies uniformly across the dense and
mixture-of-experts variants of the family.
\end{abstract}

\section{Introduction}

Adapting an open-weight base model into a capable, aligned assistant requires careful
coordination of data curation and optimization strategies across a multi-stage pipeline.
The Zen training methodology addresses the post-base lifecycle: from curated continued
pretraining through instruction tuning and alignment, to final deployment optimization.
The Zen models inherit their architecture and tokenizer from the corresponding Qwen3 base
models \cite{qwen3}; this report focuses on the data and training methodology applied on
top of those checkpoints.

Key challenges we address include:
\begin{itemize}
  \item \textbf{Data quality at scale}: Curating a continued-training corpus that balances breadth and quality
  \item \textbf{Stable optimization}: Preventing loss spikes and ensuring smooth convergence
  \item \textbf{Alignment without degradation}: Instilling safety properties while preserving capability
  \item \textbf{Compute efficiency}: Achieving good FLOPs allocation across model sizes
\end{itemize}

Because the family includes both dense variants and a mixture-of-experts variant
(Qwen3-30B-A3B, which activates a subset of experts per token), our data-mixing and
training-dynamics decisions must accommodate both dense and sparsely-activated models
under a single pipeline.

\section{Background and Related Work}

\subsection{Scaling Laws}

The Chinchilla scaling laws~\cite{hoffmann2022chinchilla} established that optimal
training requires roughly 20 tokens per parameter. Subsequent work has shown that
inference-optimal training often favors over-training smaller models on more tokens.
The Zen family follows an inference-optimal strategy, training smaller models on
significantly more data than Chinchilla-optimal prescriptions.

\subsection{Data Curation}

Prior work on data curation has emphasized the importance of deduplication, quality
filtering, and domain mixing. C4~\cite{raffel2020c4}, The Pile~\cite{gao2021pile}, and
RedPajama~\cite{together2023redpajama} established key preprocessing pipelines. Zen
extends these with perplexity-based filtering and domain-aware mixing.

\subsection{Multi-Stage Training}

The standard paradigm of pretraining followed by instruction tuning (SFT) and RLHF
was established by InstructGPT~\cite{ouyang2022rlhf}. Zen introduces constitutional
training as an intermediate stage between SFT and RLHF, reducing human annotation
requirements while improving alignment quality.

\section{Data Curation Pipeline}

\subsection{Raw Data Collection}

The Zen continued-training corpus aggregates data from multiple sources. The token
counts below describe the relative composition and the effect of the filtering pipeline;
they are illustrative of the mixture rather than a claim of from-scratch pretraining
(the base models are the pretrained Qwen3 checkpoints):

\begin{table}[H]
\centering
\begin{tabular}{lrrr}
\toprule
\textbf{Source} & \textbf{Raw Size} & \textbf{After Filter} & \textbf{Mix Weight} \\
\midrule
Web (Common Crawl) & 68T tokens & 4.2T & 60\% \\
Books \& Literature & 400B tokens & 380B & 12\% \\
Scientific Papers & 150B tokens & 140B & 8\% \\
Code Repositories & 800B tokens & 650B & 12\% \\
Wikipedia + Encyclopedias & 80B tokens & 78B & 3\% \\
Curated Q\&A & 120B tokens & 110B & 5\% \\
\midrule
\textbf{Total} & \textbf{$\sim$70T} & \textbf{7T} & \textbf{100\%} \\
\bottomrule
\end{tabular}
\caption{Zen pretraining corpus composition after filtering.}
\end{table}

\subsection{Quality Filtering}

Our quality pipeline applies filters in sequence, each removing a fraction of documents:

\begin{algorithm}[H]
\caption{Zen Quality Filtering Pipeline}
\begin{algorithmic}[1]
\REQUIRE Raw document corpus $\mathcal{D}$
\ENSURE Filtered corpus $\mathcal{D}^*$
\STATE $\mathcal{D}_1 \leftarrow \text{LanguageID}(\mathcal{D})$, keep $p(\text{en}) > 0.65$
\STATE $\mathcal{D}_2 \leftarrow \text{URLFilter}(\mathcal{D}_1)$, remove spam/adult domains
\STATE $\mathcal{D}_3 \leftarrow \text{MinHashDedup}(\mathcal{D}_2, \text{ngram}=13, \text{jaccard}>0.8)$
\STATE $\mathcal{D}_4 \leftarrow \text{HeuristicFilter}(\mathcal{D}_3)$:
\STATE \quad Remove docs $< 100$ tokens or $> 100$K tokens
\STATE \quad Remove docs with symbol/word ratio $> 0.1$
\STATE \quad Remove docs with repeated n-gram ratio $> 0.3$
\STATE Train reference LM $\mathcal{M}_\text{ref}$ on high-quality seed corpus (200B tokens)
\STATE $\mathcal{D}_5 \leftarrow \{d \in \mathcal{D}_4 : \text{PPL}_{\mathcal{M}_\text{ref}}(d) < \tau_\text{ppl}\}$
\STATE $\tau_\text{ppl} \leftarrow$ 90th percentile of $\text{PPL}_{\mathcal{M}_\text{ref}}$ on $\mathcal{D}_4$
\STATE $\mathcal{D}^* \leftarrow \text{DomainSample}(\mathcal{D}_5, \text{weights}=W)$
\RETURN $\mathcal{D}^*$
\end{algorithmic}
\end{algorithm}

\subsubsection{Perplexity-Based Quality Filtering}

A key innovation is using a reference language model to score document quality.
We train a 1.3B parameter reference model on a manually curated seed corpus of
200B high-quality tokens. Documents scoring above the 90th percentile perplexity
threshold are removed, eliminating repetitive, incoherent, or template-generated content.

The perplexity filter is highly selective on raw web content (retaining only a small
fraction of crawled tokens) while books and scientific papers show 95\%+ retention,
confirming the filter's discriminative power: it removes low-quality web text far more
aggressively than curated high-quality sources.

\subsection{Hybrid Data Mixing}

Static mixing ratios degrade performance as training progresses because the model
learns different domains at different rates. We employ adaptive mixing via a
domain loss monitor:

\begin{equation}
w_d^{(t+1)} = w_d^{(t)} \cdot \exp\left(\eta \cdot \frac{\mathcal{L}_d^{(t)} - \bar{\mathcal{L}}^{(t)}}{\sigma_\mathcal{L}^{(t)}}\right)
\end{equation}

where $w_d^{(t)}$ is the weight for domain $d$ at step $t$, $\mathcal{L}_d^{(t)}$ is
the per-domain validation loss, and $\eta = 0.01$ is the adaptation rate. Weights
are normalized after each update. This causes the optimizer to allocate more capacity
to domains where the model lags.

\section{Training Setup}

\subsection{Architecture: Dense and Mixture-of-Experts Variants}

The Zen family inherits its architecture from the Qwen3 base models \cite{qwen3}: the
0.6B/4B/8B/32B variants are dense transformers, and the 30B-A3B variant is a
mixture-of-experts (MoE) model that replaces dense FFN layers with a mixture of $E$
expert networks activated sparsely via top-$k$ routing:

\begin{equation}
\text{MoE}(x) = \sum_{i=1}^{k} g_i(x) \cdot E_i(x), \quad g(x) = \text{TopK}(\text{softmax}(W_g x), k)
\end{equation}

so that only $k$ of the $E$ experts (roughly 3B active parameters out of 30B total for
the 30B-A3B variant) are evaluated per token. This sparse activation is the source of the
MoE variant's favorable quality-per-active-FLOP, and it informs the data-mixing and
training-dynamics decisions described below.

\subsection{Model Configurations}

\begin{table}[H]
\centering
\begin{tabular}{lrl}
\toprule
\textbf{Model} & \textbf{Params} & \textbf{Type} \\
\midrule
Zen-Nano & 0.6B & dense (Qwen3-0.6B) \\
Zen-Eco  & 4B   & dense (Qwen3-4B) \\
Zen-8B   & 8B   & dense (Qwen3-8B) \\
Zen-32B  & 32B  & dense (Qwen3-32B) \\
Zen-Omni & 30B-A3B & MoE, $\sim$3B active (Qwen3-30B-A3B) \\
\bottomrule
\end{tabular}
\caption{Zen model family configurations. All variants are Qwen3-based and Apache-2.0;
detailed layer/width/head counts follow the corresponding Qwen3 base models.}
\end{table}

\subsection{Optimization}

All models are trained with AdamW ($\beta_1=0.9$, $\beta_2=0.95$, $\epsilon=10^{-8}$,
weight decay $\lambda=0.1$). Learning rate follows a warmup-cosine schedule:

\begin{equation}
\eta_t = \eta_{\max} \cdot \begin{cases} t / T_{\text{warm}} & t < T_{\text{warm}} \\ \frac{1}{2}\left(1 + \cos\left(\pi \cdot \frac{t - T_{\text{warm}}}{T - T_{\text{warm}}}\right)\right) & t \geq T_{\text{warm}} \end{cases}
\end{equation}

with $T_{\text{warm}} = 2000$ steps, $\eta_{\max}$ model-size dependent (see Table 3),
and final learning rate $\eta_{\min} = \eta_{\max}/10$.

\begin{table}[H]
\centering
\begin{tabular}{lrr}
\toprule
\textbf{Model} & \textbf{Peak LR} & \textbf{Batch Size (tokens)} \\
\midrule
Zen-Nano (0.6B) & $6\times10^{-4}$ & 2M \\
Zen-Eco (4B)    & $4\times10^{-4}$ & 4M \\
Zen-8B          & $3\times10^{-4}$ & 4M \\
Zen-32B         & $1\times10^{-4}$ & 8M \\
Zen-Omni (30B-A3B) & $1\times10^{-4}$ & 8M \\
\bottomrule
\end{tabular}
\caption{Continued-training hyperparameters per model. The MoE variant uses the same
peak learning rate as the 32B dense model but a lower effective per-parameter update due
to sparse activation.}
\end{table}

\subsection{Compute Infrastructure}

Training runs on H100 SXM5 80GB clusters with NVLink interconnects. Pipeline
parallelism depth $P$, tensor parallelism $T$, and data parallelism $D$ are
configured per model size:

\begin{table}[H]
\centering
\begin{tabular}{lrrr}
\toprule
\textbf{Model} & \textbf{P} & \textbf{T} & \textbf{D} \\
\midrule
Zen-8B           & 1 & 4 & data-parallel \\
Zen-32B          & 4 & 8 & data-parallel \\
Zen-Omni (30B-A3B) & 4 & 8 & data-parallel \\
\bottomrule
\end{tabular}
\caption{Parallelism strategies per model. The MoE variant additionally uses expert
parallelism across the 30B-A3B expert set.}
\end{table}

\section{Supervised Fine-Tuning}

\subsection{SFT Data Construction}

The SFT dataset contains 2.4M instruction-response pairs spanning:
multiturn conversation (35\%), coding tasks (25\%), mathematical reasoning (20\%),
long-form writing (12\%), and factual Q\&A (8\%). All pairs are human-validated
or model-generated with human verification.

\subsection{SFT Training Protocol}

SFT is conducted for 3 epochs with learning rate $5\times10^{-6}$, batch size 256,
and maximum sequence length 32K tokens. Only response tokens contribute to the
cross-entropy loss:

\begin{equation}
\mathcal{L}_\text{SFT} = -\sum_{t \in \text{response}} \log p_\theta(x_t \mid x_{<t})
\end{equation}

We apply NEFTune noise injection~\cite{jain2023neftune} with $\alpha=5$ to improve
generalization on open-ended tasks.

\section{Constitutional Training}

Between SFT and RLHF, we apply a constitutional training stage that teaches the model
to critique and revise its own outputs according to a set of principles:

\begin{enumerate}
  \item Generate response $r_0$ to prompt $p$
  \item Apply critique template: ``Identify ways the response violates principle $c_i$''
  \item Generate revision $r_1$ guided by the critique
  \item Train on $(p, r_1)$ pairs using SFT loss
\end{enumerate}

This substantially reduces human annotation requirements for the subsequent RLHF stage
while improving harmlessness, as the model learns to self-correct against the principles
before preference optimization.

\section{RLHF Pipeline}

\subsection{Reward Model Training}

The reward model (RM) is initialized from the Zen-8B SFT checkpoint and trained on
human preference pairs. Given two responses $r_a, r_b$ to the same prompt,
the RM is trained to maximize the margin:

\begin{equation}
\mathcal{L}_\text{RM} = -\mathbb{E}_{(p,r_w,r_l)} \left[\log \sigma\left(r_\phi(p, r_w) - r_\phi(p, r_l)\right)\right]
\end{equation}

where $r_w$ is the preferred response and $r_l$ the rejected response.

\subsection{PPO Training}

Policy optimization uses Proximal Policy Optimization with the KL penalty:

\begin{equation}
\mathcal{L}_\text{PPO} = \mathbb{E}\left[r_\phi(p, r) - \beta \cdot \text{KL}\left[\pi_\theta(\cdot|p) \| \pi_{\text{ref}}(\cdot|p)\right]\right]
\end{equation}

$\beta = 0.04$ is tuned to balance helpfulness and proximity to the SFT policy.
We apply a reward clipping of $[-5, 5]$ and run 2 PPO epochs per batch.

\section{Experiments and Results}

\subsection{Training Curves}

Continued-training loss curves exhibit three distinct phases observed across all model
sizes: (1) rapid descent early in training as the model adapts to the curated data
distribution, (2) steady decline through the bulk of training as domain knowledge
accumulates, (3) a slower final phase with diminishing returns signaling data saturation.

Loss spike prevention is achieved via gradient norm clipping at 1.0 and a loss
spike detector that halves the learning rate for 100 steps when $\mathcal{L}_t > 1.5 \cdot \bar{\mathcal{L}}_{t-100:t}$.

\subsection{Ablation Results}

\begin{table}[H]
\centering
\begin{tabular}{lc}
\toprule
\textbf{Configuration removed from full pipeline} & \textbf{Capability drop (MMLU/HumanEval/MATH/MT-Bench)} \\
\midrule
-- (full pipeline)         & reference (best) \\
No PPL filter              & largest drop \\
Static mix (no adaptive)   & moderate drop \\
No constitutional training & small drop \\
No NEFTune                 & small drop \\
\bottomrule
\end{tabular}
\caption{Component ablation on Zen-8B (Qwen3-8B): effect of removing each methodological
component from the full pipeline, ordered by the size of the resulting capability drop
across MMLU, HumanEval, MATH, and MT-Bench. The perplexity filter is the single most
important component.}
\end{table}

\subsection{Compute Budget Analysis}

We analyze the continued-training compute-to-performance tradeoff by fitting a saturating
power law:

\begin{equation}
\text{Performance}(C) \approx A - B \cdot C^{-\alpha}
\end{equation}

where $C$ is the training compute in FLOPs. Empirically, downstream MMLU performance
saturates with additional continued-training compute for a fixed model size, exhibiting
clear diminishing returns; reaching higher capability targets requires moving to a larger
model in the family rather than spending more compute on a small one.

\section{Analysis}

\subsection{Data Quality vs. Quantity}

Our experiments confirm that data quality dominates quantity beyond a threshold:
\emph{halving} the corpus size with a tighter perplexity filter
($\tau_\text{ppl}$ at the 80th percentile) yields a larger MMLU improvement than
\emph{doubling} the corpus size with the filter disabled. Aggressive quality filtering
is therefore preferable to indiscriminate scale.

\subsection{Expert Specialization in the MoE Variant}

Analysis of expert activation patterns in the Zen-Omni (Qwen3-30B-A3B) MoE variant
reveals domain specialization: code-related tokens activate a consistent subset of
experts, mathematics activates a partially overlapping set, and natural language
distributes more broadly. This specialization emerges from the data distribution during
continued training without explicit routing supervision.

\subsection{Scaling Behavior}

Zen models follow a modified scaling law that accounts for MoE efficiency:

\begin{equation}
\mathcal{L}(N_\text{active}, D) = \frac{A}{N_\text{active}^\alpha} + \frac{B}{D^\beta}
\end{equation}

where $N_\text{active}$ is the number of active parameters per token (not total
parameters) and $D$ is the dataset size. Cast in terms of active parameters, the MoE
variant achieves lower loss at fixed active-parameter compute than a dense model of the
same active size, which is the efficiency rationale for including the 30B-A3B variant in
the family.

\section{Conclusion}

The Zen training methodology demonstrates that careful data curation, adaptive mixing,
constitutional training, and staged alignment, applied on top of the open-weight Qwen3
base models, produce models that are both capable and safe. The perplexity-based
filtering provides the largest single capability gain among the methodological components
we study. Constitutional training between SFT and RLHF reduces annotation costs while
improving alignment quality, making the pipeline more scalable.

Future work will explore continual pretraining strategies, domain-specific fine-tuning
at scale, and automated data quality assessment to further reduce human annotation
requirements in the training pipeline.

\bibliographystyle{plain}
\begin{thebibliography}{99}
\bibitem{hoffmann2022chinchilla} Hoffmann et al. (2022). Training Compute-Optimal Large Language Models. \textit{NeurIPS}.
\bibitem{raffel2020c4} Raffel et al. (2020). Exploring the Limits of Transfer Learning with a Unified Text-to-Text Transformer. \textit{JMLR}.
\bibitem{gao2021pile} Gao et al. (2021). The Pile: An 800GB Dataset of Diverse Text for Language Modeling. \textit{arXiv:2101.00027}.
\bibitem{together2023redpajama} Together AI (2023). RedPajama: An Open Source Recipe to Reproduce LLaMA Training Dataset. \textit{GitHub}.
\bibitem{ouyang2022rlhf} Ouyang et al. (2022). Training language models to follow instructions with human feedback. \textit{NeurIPS}.
\bibitem{jain2023neftune} Jain et al. (2023). NEFTune: Noisy Embeddings Improve Instruction Finetuning. \textit{arXiv:2310.05914}.
\bibitem{qwen3} Qwen Team (2025). Qwen3 Technical Report. \textit{arXiv:2505.09388}.
\end{thebibliography}

\end{document}
